\section{I Sleep, but the Heart is Awake}

\begin{quotex}
I sleep, but my heart is awake. 
\flright{\textit{Song of Solomon, 5:2}}
\end{quotex}

Please consider the four states of consciousness mentioned by \textbf{Bede Griffiths} in \textit{The Cosmic Mystery}\footnote{\url{https://gornahoor.net/?p=8055}}.


We are accustomed to consider the sleeping, dreaming, and waking states as distinct states of consciousness, which follow each other consecutively. However, most of us have noticed that, rather than sequential, they interpenetrate each other. Our presumed waking state is mostly a series of semi-dreams: i.e., thoughts, imaginations, little stories, and so on. Then there is something deeper, beyond conscious awareness, in operation, affecting our postures, movements, attention and so on.

Among those dreams is the dream of awakening. That ``I'', the subject of that dream, will ``snap us out of it'' from time to time. That I needs to be nourished until the whole heart is awake.

\paragraph{Theology of the Body}


A delay was caused by some new information on the \textit{Theology of the Body}\footnote{\url{https://sspx.org/en/theology-of-the-body-genesis-of-confusion}} (H/T Perennial) that I came across as I was writing the introduction to the chapters on Sintesi that relate the sexes to race. Surprisingly, Evola's ideas have some things in common with \textbf{Pope John Paul II}'s conceptions. Perhaps that is so, given the influence of German idealism and the philosophy of personalism on both of them.

For Scholasticism, the mind, or intellectual soul, is the image of God. The mind is without qualities, neither male nor female. For JPII and Evola, on the other hand, the body is the outward expression of spirit, hence either masculinity or femininity permeates the whole being. The Scholastic approach easily devolves to a Cartesian dualism of body and mind.

Moreover, the Scholastic definition of the person is too abstract: \emph{an individual substance of a rational nature}. It does not mention Will, I, Freedom, Consciousness. Then there is Guenon's objection that it focuses on Being to the exclusion of non-Being. Specifically, it is too static, not making the distinction between the virtual and the actual. The actualization of a person is something else. Will must be developed, consciousness expanded, freedom/liberation must be achieved as it is not simply a given. ``Truth shall set you free.'' For the Scholasitics, there can be no development of the person. However, for JPII and Evola one of hallmarks of the person is self-possession, which requires a development.

Also, speculative theology does not make a clear distinction between soul and spirit. Mystical theology, on the other hand, does. \textbf{St John of the Cross}, for example, distinguishes between psychic experiences and truly spiritual experience. \textbf{St John of the Cross} is one of the influences on the \textit{Theology of the Body}.

\paragraph{Sex and the Body}


These are some of the main points of the Theology of the Body. When the Evola translations come out, we can see better how they relate.

\begin{itemize}


\item There is no abstract ``human'': the human being is always embodied as a man or a woman. Masculinity and Femininity are expressions of the person.

\item The Spirit is not without qualities.

\item The body is not just matter, subject to physical or biological laws. Rather, it is permeated by the soul and the spirit. Otherwise, resurrection cannot be understood except as the revival of a corpse.

\item Biology is a human construct and does not exhaust all we know about the human being, sex, or the body.

\item The Baconian project of gaining power over nature, specifically, technological power over the body, is rejected. Rather, the subtle rules the dense so the spirit must control the body, without relying on material means.
\end{itemize}


\paragraph{Unintended Consequences}


There are perhaps other results that can be derived from this approach, not that they can be attributed to JPII. Biblical, mystical, philosophical, and phenomenological sources are used to develop the theology of the body. Why, then, can we not include other sources and even traditional sources? For example, Evola relies on \textbf{Otto Weininger}'s book \textit{Sex and Character} as well as the Laws of Manu for his own teachings on the spiritual and bodily aspects of masculinity and femininity.

Although JPII does not mention this, if the spirit has qualities, then race, or \emph{ethnos} as \textbf{E Michael Jones} calls it, has its counterpart in the spirit. The relationship is not necessarily one-to-one, as we will see. Of course, this doctrine rejects the idea that biology or DNA is the cause of differences in psychic and spiritual makeup of the person.

Furthermore, the body cannot be understood as a piece of matter, but rather as a non-dual body/soul/spirit complex. This conception is developed by \textbf{Valentin Tomberg} in Letter XX of the Meditations of the Tarot. He shows the chain:

spirit $\Rightarrow$ psychic forces $\Rightarrow$ energy $\Rightarrow$ material organs.

That is the vertical process that works with the horizontal process of heredity. More research needs to be done to relate these various currents of thought.

Then, too, the notion that the body can be more or less dense begins to make some sense. Hence, there may be no physical traces of the Hyperboreans for the reason that the body in that era was less dense. This is something proposed by Rene Guenon. In this case, the positive sciences cannot be the last word.


\flrightit{Posted on 2015-04-21 by Cologero}

\begin{center}* * *\end{center}

\begin{footnotesize}
\begin{sffamily}
\texttt{X on 2015-04-22 at 01:16 said:}

``St John of the Cross, for example, distinguishes between psychic experiences and truly spiritual experience.''

Can you please say it more detailedly for this?

\hfill

\texttt{Cologero on 2015-04-22 at 05:36 said:}

X, this topic is discussed at length at
\textit{Three Ages of the Interior Life}\footnote{\url{http://www.christianperfection.info/index.php}} by Fr Garrigou-Lagrange.

For example, concerning psychic phenomena: \textit{The Illuminative Way of Proficients}\footnote{\url{http://www.christianperfection.info/tta82.php}}.

Also, on mystical union, which can also be read as the antidote to Cassiodorus' objections: \textit{The Transforming Union}\footnote{\url{http://www.christianperfection.info/tta105.php\#bk2}}

\hfill

\texttt{X on 2015-04-22 at 06:00 said:}

Thanks you very much, Mr. Salvo.

And the by the way, there is error in the first link.

\hfill

\texttt{X on 2015-04-22 at 06:04 said:}

``The state which St. John of the Cross describes here is a state of love linked to a state of infused contemplation. The connection is owing to a necessity of love: `True and full love cannot hide anything.' This connection is not accidental, since this need is connatural to perfect charity.''

\hfill

\texttt{thomas walker on 2015-04-22 at 10:11 said:}

We know from our own experience that body and soul are intimately related : bodily ailments can completely dominate our thoughts and emotions . I think that the kabbalistic idea of nephesh speaks to this , Our personalities are intimately linked with our bodies . There is a problem to understand how after death only the soul will exist until it is joined to the body in the Last Judgement . If time and space are abolished how can there be bodies which as bodies with extension must occupy space ?

Also , doesn't Jesus teach that in heaven there will neither be taking or giving in marriage because bodies will be like the angels ( i assume that means without gender ) Therefore gender is only accidental not essential .

\hfill

\texttt{obscure on 2015-04-22 at 15:20 said:}

Nous is indeed prior to qualities, but qualities are received prior to the quantitative features of the body since quantity most signifies prime matter. The first relation to the body and thus to the macrocosm is qualitative, but pure Nous is not for this reason qualitative. Prime matter is also not seperable and thus not really distinct from information. Although form itself as Spirit or Idea is really distinct from matter or informed matter.

The rational soul is a discursive intellect animating a material substance or complex subject (these terms are synonymous). Its constitutive features are active intellect, passive intellect and sense-perception. Pure intellect is only an active intellect and its constitutive features are its simple form and its essential act (sometimes erronesouly translated by neo-Thomists in the tradition of Gilson as `act of existence', but in truth `essential act' or `actus essendi' is pure will). A pure intellect is necessarily passive in relation to the First, but this is due to metaphysical necessity rather than discursive intellection with respect to the agent. The psychological (or `epistemological') triad of the rational soul is not to be confused with the physiological triad (rational soul, sentient soul and vegetative soul) or the moral-constitutive triad (intellect, will and appetite).

Substantial form is the subsisting simple subjectivity of a complex subject (or compound substance, material substance, etc.: These terms are synonymous). A per se simple subject (not subsisting in a complex subject) is either a Spirit or Idea which is specifically different from the Absolute simple subject. Thus the metaphysical triad of all reality is Absolute simple subject, relative or `specific' simple subjects and lastly the complex subjects. In the order of pure simple subjects Spirits are active while Ideas are passive, although with respect to what is inferior (the order of complex subjects, the material cosmos, temporal procession, nature, etc.: These terms are synonymous) Ideas are obviously called `active', but this is due to metaphysical necessity (which is to say, the hierarchy between God, supernature and nature). Corporeal form is constituted by extension and quality; it is an accidental form of a simple subject and it is the basic information of matter. Substantial form is not corporeal form, it `contains' essential possibilities not corporeal matter.

The basic qualification of Spirits other than the Absolute is their essential act or will gratuitously caused by the First (for the only relationship between the Infinite and the finite is gratuitous i.e. grace). The Absolute cannot be distinguished from His essential act in any respect since the Absolute is a Pure Act of Infinite Power. Even though neo-Thomists may not write in this way it is all to be found in the primary texts of their Doctor and of other medieval doctors. Duns Scotus is in many respects the master and I recommend this blog containing some historical materials if any are interested:

\url{http://lyfaber.blogspot.com/}


\hfill

\texttt{David on 2015-04-22 at 16:30 said:}

Cologero : Would it be possible to have all the texts within this site on a .zip or .rar format have a download link ? If the site is to end, a great part of the first texts I never had the chance to read. I imagine you have either a backup of .html files or the actual text files somewhere which could be given. Thank you.

\hfill

\texttt{obscure on 2015-04-22 at 18:48 said:}

Here are some translations of the Scotistic Theoremata:

\url{http://lyfaber.blogspot.com/2010/01/theoremata-scoti-partes-i-v.html}


Scotus' discusson of quality is interesting, although I'm sure it won't surprise many.

\hfill

\texttt{Cologero on 2015-04-22 at 20:28 said:}

TW, this is the full quote: ``In the resurrection they take neither wife nor husband, but are like the angels in heaven.''

I suppose one approach is Guenon's: the being that was in the human state is now in an angelic state.

But. JPII contrasts the resurrection body with our bodies as they are now, with its members at war with the spirit. The resurrection body will not experience that opposition, it will be spiritualized. JPII is closer to Tomberg with this explanation:

\begin{quotex}\footnotesize

Spiritualization signifies not only that the spirit will master the body, but that it will also fully permeate the body and the powers of the spirit will permeate the energies of the body. (67:1)
\end{quotex}


He takes pain to emphasize that this state is not ``disincarnation'' of the body, nor a dehumanization. There is a deep harmony between spirit and body, but primacy resides in the spirit. He goes on to relate this to theosis (translated as divinization).

In the resurrection, the body is outside of human history ``tied to marriage and procreation''. However, this does not mean that the body is suddenly neutered. (It would be interesting to compare that idea with the way Mouravieff includes marriage and procreation as part of the ``General Law''.)

\hfill

\texttt{Cologero on 2015-04-22 at 20:47 said:}

Thanks for that synopsis, Obscure, but the issue is the whether or not JPII is introducing novelty into that doctrine. It appears that he is rejecting the notion that the Spirit is distinct from the body. For the Scotist, is masculinity or femininity an accidental or an essential quality? If God has qualities, then why would not a man have them?

As for your other link, I read the claim that the ``more perfect man has more perfect intelligence''. I thought the Scotist position is that the goal of creation is the divinization of man, i.e., union with God. I don't see them as the same thing.

\hfill

\texttt{Cologero on 2015-04-23 at 08:55 said:}

David, Gornahoor will not be taken off-line, since it is intended for a future generation that may have a better perspective. However, I won't be updating it but perhaps someone else will.

\hfill

\texttt{thomas walker on 2015-04-23 at 13:14 said:}

Thanks for your reply Cologero

\hfill

\texttt{obscure on 2015-04-23 at 23:15 said:}

Hello Cologero,

Keep in mind that when I write some long comment I intend a general audience and not just yourself in particular, for I enjoy sharing information. There is much ink I could spill with regard to the hierarchies of metaphysical qualification. This would require a fine treatment with some more effort than I can provide at the moment. If anything I've provided has produced some degree of confusion it was due to a degree of privation on my part. I promise to address your points at a later time.

As for the ToB, I can't say I've read it. My impression was that it was mostly an exoteric text of a more-or-less existential nature.

I will leave off for now with a Hermetic passage (Although it is not the sort of passage to be taken with little consideration or lack of interpretation):

``Soul enters the body by necessity, Intellect enters soul by judgment. While being outside of the body, soul has neither quality nor quantity. Once it is in the body it receives, as an accident, quality and quantity as well as good and evil for matter brings about such things.''

\hfill

\texttt{Cologero on 2015-04-24 at 00:07 said:}

No confusion at all, Obscure, what you describe is how I would have been accustomed to look at it.

However, in light of the ongoing translation of Evola's Sintesi, there are questions that arise. I am not trying to provide the basis for a sophisticated theory of ``white nationalism'', but in exploring the metaphysical ramifications of his doctrine. Precisely, is there spiritual differentiation among men?

It is not so obvious that corporeal life is an accident. For example, in Plato's Myth of Er, the circumstances of birth are determined by the prior state of the spirit. For Guenon, each human being is a possibility of manifestation in God's mind, and this determines the ``accidents'' of birth. Perhaps these are differences in terms, but maybe not.

Based of the way it's been used, I, too, assumed that ToB was sort of a Catholic Kama Sutra. However, it is actually heavily indebted to John of the Cross. It also brings back final and formal cause to the understanding of the human being, contra both science and the entire modern project. It restores the notion of interiority to the body, which therefore cannot be treated simply as passive matter.

\hfill

\texttt{Michael on 2020-04-21 at 08:01 said:}

``Biology is a human construct and does not exhaust all we know about the human being, sex, or the body.''

An interesting point when contrasted to the modern approach of ``constructs''. While the modern approach says we can determine these for ourselves, JPII says there are higher methods of understanding to more great knowledge of the categories, a theology of the body (unified).

Yet in the other approach, we believe we can determine that human construct by just thinking it, isn't this just intellectual thinking centered from the individual? Pulling it down into discursive thought where all differences are forgotten instead of realized.

\end{sffamily}
\end{footnotesize}

